% !PRINTED 74
\weberpage{74}

%<<CONT
worin

\[
\varepsilon_0 = e^{\frac{\pi i \omega}{4}}
\]

und

\weq{11}{
\vartheta_{11}(0) = 0,\quad \vartheta_{10}\!\left(\tfrac{1}{2}\right) = 0,\quad \vartheta_{01}\!\left(\tfrac{\omega}{2}\right) = 0,\quad \vartheta_{00}\!\left(\tfrac{1+\omega}{2}\right) = 0.
}

\gesperrt{Die Quadrate der Funktionen (7):}

\weq{12}{
\vartheta_{00}^2(v),\quad \vartheta_{01}^2(v),\quad \vartheta_{10}^2(v),\quad \vartheta_{11}^2(v)
}

sind $\Theta_{00}$-Funktionen zweiter Ordnung, und folglich bestehen
zwischen ihnen zwei lineare Relationen. Nehmen wir diese Relationen
 in der Form an:

\[
\begin{aligned}
\vartheta_{10}^2(v) &= A\,\vartheta_{01}^2(v) + B\,\vartheta_{11}^2(v),\\
\vartheta_{00}^2(v) &= A'\vartheta_{01}^2(v) + B'\vartheta_{11}^2(v),
\end{aligned}
\]

so kann man die Koeffizienten auf Grund der Formeln (10), (11)
leicht bestimmen, wenn man $v = 0$ und $v = \dfrac{\omega}{2}$ setzt. Man
erhält so

\weq{13}{
\begin{aligned}
\vartheta_{01}^2\,\vartheta_{10}^2(v) &= \vartheta_{10}^2\,\vartheta_{01}^2(v) - \vartheta_{00}^2\,\vartheta_{11}^2(v),\\
\vartheta_{01}^2\,\vartheta_{00}^2(v) &= \vartheta_{00}^2\,\vartheta_{01}^2(v) - \vartheta_{10}^2\,\vartheta_{11}^2(v),
\end{aligned}
}

und daraus noch, indem man $v = \tfrac{1}{2}$ setzt:

\weq{14}{
\vartheta_{00}^4 = \vartheta_{01}^4 + \vartheta_{10}^4.
}

Durch zwei beliebige von den Quadraten (12) können alle
Funktionen $\Theta_{00}$ der zweiten Ordnung linear dargestellt werden;
aber auch die übrigen $\Theta$-Funktionen zweiter Ordnung lassen sich
aus den Funktionen $\vartheta_{g_1,g_2}$ zusammensetzen, denn man erhält für
jede Charakteristik zwei linear unabhängige Produkte, von denen
das eine eine gerade, das andere eine ungerade Funktion ist,
nämlich:

\weq{15}{
\begin{aligned}
\vartheta_{00}(v)\,\vartheta_{01}(v), &\quad \vartheta_{10}(v)\,\vartheta_{11}(v) & &\text{Charakteristik } (0,1)\\
\vartheta_{00}(v)\,\vartheta_{10}(v), &\quad \vartheta_{01}(v)\,\vartheta_{11}(v) & &\text{„}\quad (1,0)\\
\vartheta_{10}(v)\,\vartheta_{01}(v), &\quad \vartheta_{00}(v)\,\vartheta_{11}(v) & &\text{„}\quad (1,1)
\end{aligned}
}

Nach demselben Prinzip lassen sich nun alle $\Theta$-Funktionen
beliebiger Ordnung und beliebiger Hauptcharakteristik aus den
Funktionen $\vartheta$ bilden. Um dies nachzuweisen, bezeichnen wir mit
$\Theta_0, \Theta_1$ irgend zwei der $\vartheta$-Quadrate (12) und mit $F^{(n)}(\Theta_0,\Theta_1)$ eine
ganze rationale und homogene Funktion $n$ter Ordnung der beiden
Argumente $\Theta_0, \Theta_1$. Man erhält hiernach für die $\Theta$-Funktionen
$m$ter Ordnung $\Theta^{(m)}(v)$ folgende Ausdrücke:

%CONT>>
